\section{An independent gradient argument below one half}
\label{sec:rcut-introduction}
These appendices prove cutoff for $0\le\beta<1/2$ by following the
unweighted gradient along the heat-bath dynamics. The proof places the
median mixing time on the logarithmic scale and shows that the
total-variation transition has vanishing relative width. The central
estimate is pointwise: the gradient dissipates at a fixed positive rate outside a set
of exponentially small Gibbs mass. Entropy decay first makes that
exceptional set unlikely, and smoothing over a positive logarithmic lag
then converts gradient control into the cutoff ratio.

This argument also yields estimates with uses beyond its cutoff
conclusion. They include a spectral-gap and covariance bound uniform over
external fields and interaction scalings, and the scalar variational
bound $0.9969$ with its complete analytic and arithmetic certificate.
The proof is independent of the limiting exact-gradient operator used in
the full high-temperature argument.

Keep the model, half-differences, and two clock conventions of
Section~\ref{sec:model}. Write $\P_J$ for the interaction-matrix law and
abbreviate the worst-start total-variation distances by
$d_c(t)=d_n(t)$ and $d_d(k)=d_n^{\rm disc}(k)$. We will use the two
gradient densities and the median mixing time
\[
 U(f)=\sum_i(d_i f)^2,\qquad
 \Gamma(f)=\sum_iw_i(d_i f)^2,\qquad
 t_n^{\rm med}=\inf\{t\ge0:d_c(t)\le1/2\}.
\]
The weights are $w_i=1-x_i\tanh((Jx)_i)$. Expanding the square at each
spin flip gives the identity
\begin{equation}
 L(f^2)-2fLf=2\Gamma(f)
 \label{rcut:eq:variance-density}
\end{equation}
For a probability measure $\nu$ and an integrable function $f$, write
$\nu f=\int f\,d\nu$.
All vector and matrix norms in these appendices are Euclidean and
Euclidean operator norms, respectively, unless explicitly specified.
We continue to write $\P_x$ and $\E_x$ for chain probabilities and
expectations after fixing the disorder.

\begin{theorem}\label{rcut:thm:main}
For every fixed $\beta\in[0,1/2)$, $\varepsilon\in(0,1/2)$, and
$\eta>0$, let
\[
 t_{n,\beta,J}(\varepsilon)=\min\{k\ge0:d_d(k)\le\varepsilon\}.
\]
Then
\[
 \P_J\left\{
 \frac{t_{n,\beta,J}(\varepsilon)}
 {t_{n,\beta,J}(1-\varepsilon)}>1+\eta\right\}\longrightarrow0.
\]
The denominator is positive for sufficiently large $n$.
\end{theorem}

The denominator is eventually positive because
$d_d(0)=1-\min_x\mu(x)\ge1-2^{-n}$. We prove the stated ratio directly;
the full-phase theorem is not an input to this appendix argument.

The proof has three stages. Section~\ref{rcut:sec:entropy} starts with a
uniform two-spin inequality, obtains the gap under all external fields
and interaction scalings, and passes from that gap to entropy decay by
Gaussian observations. Section~\ref{rcut:sec:gradient} expresses the
pointwise gradient dissipation as a matrix inequality. The matrix
comparison in Section~\ref{rcut:sec:comparison} and the scalar estimate in
Section~\ref{rcut:sec:scalar} establish that inequality outside an
exponentially small Gibbs set. Finally,
Sections~\ref{rcut:sec:dynamics} and~\ref{rcut:sec:cutoff} use concentration
and positive-lag smoothing to prove the cutoff ratio, invoking the
general jump-moment and clock lemmas stated in Section~\ref{sec:model}.
Section~\ref{rcut:sec:scalar-certificate} supplies every quadrature and
arithmetic enclosure used in the scalar estimate.

The temperature parameter $\beta$ remains fixed throughout, and constants
may depend on it. Whenever an estimate is used on a logarithmic time
horizon, its times are at most $D_*\log n$, with $D_*<\infty$ chosen
before taking $n\to\infty$.
